\documentclass[book.tex]{subfiles}
\begin{document}


\chapter{Numerical Solution of Partial Differential Equations}

\label{chap:numerical_pde}
\noindent\rule{11cm}{0.4pt} \\
\textbf{Prerequisites:} \autoref{chap:pde} \\
\textbf{Difficulty Level:} ** \\
\noindent\rule{11cm}{0.4pt}
\vspace{5mm}

In this chapter, we briefly review techniques for solving partial differential equations numerically.

\section{Finite Element Methods}

Finite element methods divide up the solution space into a finite mesh. Mathematically, a mesh is a collection of points
\begin{align}
\mathcal{M} &= \{ (x_i,y_i,z_i,\dotsc) \}
\end{align}
which closely tile an input shape or space. Mesh elements are typically connected to their nearest neighbors to induce a graph structure. A finite element method can be rougly broken down into the recipe

\begin{enumerate}
    \item Generate mesh for domain
    \item Do iterated updates of discretized versions of the original PDE equations on the mesh
\end{enumerate}

A broad range of methods exists to compute mesh discretizations of suitable shapes. These methods can work well in lower dimensional problems but can struggle at high dimensions.

\section{Mesh Free Methods}

A newer class of numerical methods to solve PDEs use mesh-free representations for solutions \cite{fasshauer2006meshfree}. Simple mesh-free methods represent the value of a solution globally by reference to a set of control points

\begin{align}
u(x,t) &= \sum_i m_i u_i W(x, (x_i,\dotsc))
\end{align}
where $u_i$ is the value at the $i$-th control point, and $W$ is a distance kernel function. This representation is adapted by many differentiable PDE solver approaches.




\end{document}